\section{Contribution Taxonomy}

The scaffold makes seven distinct contributions to machine-learning research practice. First, it provides a scientific-process contribution by re-centering machine-learning work on the sequence from question to design, evidence, claim, and publication. Second, it provides a reproducibility contribution by treating each stage of the research process as a versioned artifact rather than as an informal project note. Third, it provides a claim-discipline contribution by requiring manuscript claims to be linked to evidence sources and bounded by the strength of validation. Fourth, it provides a reporting-operationalization contribution by translating reporting expectations into upstream workflow steps. Fifth, it provides a manuscript-engineering contribution by treating the manuscript as a structured research output generated from prior artifacts. Sixth, it provides a pedagogical contribution by making the hidden structure of publishable machine-learning research visible to students and applied researchers. Seventh, it provides a review-support contribution by giving collaborators, supervisors, reviewers, and editors a traceable path from design decisions to final claims.

\begin{table}[htbp]
\centering
\caption{Contribution taxonomy for the query-to-publication scaffold}
\begin{adjustbox}{max width=\textwidth}
\begin{tabularx}{\textwidth}{p{0.24\textwidth}X X}
\toprule
\textbf{Contribution type} & \textbf{Problem addressed} & \textbf{Scaffold response} \\
\midrule
Scientific process & Machine-learning projects often begin with models rather than research questions. & The scaffold begins with query framing, study protocol, and claim boundaries. \\
Reproducibility & Reproducibility materials are often incomplete or detached from the manuscript. & Each workflow stage produces a versioned artifact that can be reviewed and preserved. \\
Claim discipline & Manuscript claims may exceed the evidence produced by validation. & Claims are assigned levels and linked to supporting artifacts and boundary notes. \\
Reporting operationalization & Reporting standards describe final disclosures but not always how to produce them. & The scaffold turns reporting expectations into upstream workflow stages. \\
Manuscript engineering & Manuscripts are often written after analysis from loosely organized notes. & Manuscript sections are built from structured research artifacts. \\
Pedagogy & Responsible machine-learning research workflows are often implicit. & The scaffold makes the process teachable and inspectable. \\
Review support & Reviewers may struggle to trace claims back to design and evidence. & The scaffold creates a visible trail from research decisions to manuscript language. \\
\bottomrule
\end{tabularx}
\end{adjustbox}
\end{table}

The resulting contribution is workflow-level rather than algorithmic. The scaffold does not propose a new estimator or validation metric. Instead, it operationalizes reproducible machine-learning science as a query-to-publication workflow in which each major claim is traceable to a staged artifact and bounded by the strength of its supporting evidence.
